\section{Related Work}

\subsection{3D Human and Hand Motion Forecasting}

Forecasting articulated motion has been studied extensively in the broader human-motion literature. Early and influential deterministic models relied on recurrent or feed-forward temporal prediction \cite{Martinez2017RNN}, while later work improved long-range structure through motion attention \cite{Mao2020HistoryRepeats}, graph reasoning \cite{Dang2021MSRGCN}, and stronger feed-forward sequence modeling \cite{Guo2023SiMLPe}. In parallel, stochastic methods were introduced to represent the multi-modal nature of future motion through latent-variable modeling \cite{Yuan2020DLow,aliakbarian2020stochastic}, diverse-hypothesis generation \cite{mao2021gsps,Ma2022MultiObjective}, and probabilistic multi-modal forecasting \cite{Salzmann2022Motron}. These directions establish the main trade-off that also appears in hand forecasting: deterministic models often capture dominant trends well, whereas stochastic models are better suited to ambiguous futures.

Hand forecasting narrows this problem to a more densely articulated and manipulation-oriented signal. Compared with whole-body forecasting, hand motion contains stronger local couplings among nearby joints and sharper deformations associated with grasping or interaction. Recent egocentric hand-trajectory works study future hand motion as an anticipatory signal for interaction understanding \cite{Liu2022Hotspots}, egocentric 3D hand trajectory forecasting \cite{Bao2023USST,Hatano2025EgoH4}, diffusion-based hand prediction \cite{Tang2024PromptFDDM}, and interaction-aware trajectory reasoning \cite{Chen2025EgoMAN}. This paper considers that hand-specific setting and models future 3D hand motion through a deterministic coarse trend followed by stochastic residual refinement.

\subsection{Frequency-Domain Motion Modeling}

Frequency-domain representations provide a compact way to model smooth motion trends and long-range temporal structure. In particular, DCT-based forecasting has shown that a trajectory can often be predicted effectively from a low-dimensional spectral representation, reducing the burden on the predictor to model frame-level variation directly \cite{Guo2023SiMLPe}. This idea is attractive for hand forecasting because a substantial portion of the predictable trajectory trend can be captured by low-frequency components, while higher-frequency variation is harder to model deterministically. Our method builds on this perspective by using a frequency-domain coarse predictor. Unlike methods that use frequency representations only as a compact deterministic forecasting space \cite{Guo2023SiMLPe}, \modelname{} also defines residual stochastic refinement over future-frequency coefficients.

\subsection{Diffusion Models for Motion Forecasting}

Diffusion models provide a flexible framework for modeling structured uncertainty through iterative denoising \cite{Ho2020DDPM,Song2021DDIM,Salimans2022ProgressiveDistillation}. They have been used for motion generation \cite{Tevet2023MDM}, motion completion for prediction \cite{Chen2023HumanMAC}, stochastic human motion forecasting \cite{Sun2024CoMusion}, behavior-driven motion prediction \cite{barquero2023belfusion}, and nonisotropic diffusion for 3D motion prediction \cite{Curreli2025SkeletonDiffusion}. For forecasting, diffusion is especially appealing because it can generate multiple plausible futures conditioned on the observed past. However, these formulations are commonly introduced without an explicit hand-structured spatio-spectral prior. \modelname{} differs by centering diffusion around a deterministic coarse forecast and by structuring the residual noise process through hand topology and future-frequency modes.

\subsection{Hand-Centric Anticipation in Egocentric and Interactive Settings}

Egocentric and interactive perception settings motivate forecasting models that anticipate how people will move before an action is complete. In egocentric vision, hand trajectories provide a direct cue about future interaction with objects and the surrounding environment \cite{Bao2023USST,Chen2025EgoMAN}. In human--robot collaboration, related work treats probabilistic motion prediction as a component for anticipatory assistance and safer interaction \cite{Kothari2023ConstrainedHRCMotionPrediction}. In teleoperation, uncertainty-aware intention prediction serves a related anticipatory role \cite{Heman2026TeleoperationIntent}. These settings motivate hand-motion forecasting models that are both stochastic and structurally constrained. Our work remains focused on 3D hand-motion forecasting, but adopts this anticipatory perspective by modeling future hand motion as a distribution over articulated residual trajectories.
